\par\bigskip\Needspace{8\baselineskip}
\pdfbookmark[1]{Störungen der Ceres: Bemerkungen}{ceres-bemerkungen}
\SourceLeaf{416}
\begin{center}BEMERKUNGEN.\end{center}
\begingroup\small
\noindent\hspace*{\parindent}Die Notiz [2] des Abschnitts I stammt aus einem Heftchen mit dem Titel Schedae, die Notizen [3]---[7] aus einem Handbuch. Wegen der fragmentarischen Beschaffenheit der Originale musste die Bearbeitung recht frei ausfallen. Das genannte Handbuch enthält auch die erste numerische Rechnung der Störungen nach den in [7] angegebenen Formeln. Ihre Ergebnisse sind in \textsc{Zachs} Monatlicher Correspondenz, Bd. VI (Werke, Bd. VI, S. 225), abgedruckt. Auf der Grundlage dieser Störungen hat \textsc{Gauss} die Elemente erneut verbessert (Bd. VI, S. 228 und Notiz [1]) und dann mit den verbesserten Elementen die Störungen neu berechnet. Diese letzte Rechnung findet sich auf einigen zusammengehefteten Blättern mit dem Titel »Praktische Anweisung zur Berechnung der Störungen«. Sie enthalten eingangs auch eine Zusammenstellung der verwendeten Formeln einschließlich derjenigen für die Breitenstörungen (Notiz [8]). Diese hatte \textsc{Gauss} in der ersten Rechnung beiseitegelassen. Ihre Herleitung wurde nicht gefunden. Die Ergebnisse der zweiten Rechnung sind in \textsc{Zachs} Monatlicher Correspondenz, Bd. VI (Werke, Bd. VI, S. 228---229), abgedruckt. \textsc{Gauss} setzte die Jupitermasse gleich \(\dfrac{1}{1067{,}09}\) (nach \textsc{Laplace}).

Die unter II abgedruckten Briefstellen beziehen sich auf die Tafeln, die \textsc{Gauss} im Anschluss an die genannten Rechnungen erstellt und in \textsc{Zachs} Monatlicher Correspondenz, Bd. VII (Werke, Bd. VI, S. 235---243), veröffentlicht hat.

Zwischen den beiden Briefen an \textsc{Olbers} vom 25. Januar und 10. Mai 1805 (oder früher) begann \textsc{Gauss}, seine erste Methode auf die zweiten Potenzen der Exzentrizitäten auszudehnen. Der Beginn dieser Entwicklungen findet sich ebenfalls im genannten Handbuch. Wegen ihres Umfangs gab \textsc{Gauss} sie jedoch sehr bald wieder auf. Darauf folgte die Rechnung nach der zweiten, unter III mitgeteilten Methode. Vgl. hierzu den Brief an \textsc{Olbers} vom 10. Mai 1805.

Bei dieser zweiten Methode wandte \textsc{Gauss} bereits die spätere \textsc{Hansen}sche Entwicklung der Störungsfunktion durch Interpolation an (vgl. \textsc{Gauss} an \textsc{Hansen}, 11. März 1843). Die Methode bildet auch die Grundlage der späteren Rechnung der Pallasstörungen. Von ihr ließen sich im selben Handbuch nur die numerischen Rechnungen finden, anhand derer die Entwicklungen rekonstruiert wurden. Da die Notizen [2]---[5] somit ganz neu verfasst wurden, sind sie in Petitsatz gedruckt. Aus dem Original stammen nur auf S. 402 die Werte von \(N\), \(\sin J\), \(\AscendingNode\) und die Erklärung der Bedeutung von \(\alpha,\epsilon,\epsilon'\), auf S. 403 die Gleichung 2) und der Wert der Konstanten bei Gleichung 3) sowie auf S. 405---407 die numerischen Werte mit Ausnahme derer auf S. 405, Zeile 9---5 von unten, und der Faktoren von \(\cos(3D+\text{\AstroSymbols ⚳}-\AscendingNode+M)\) auf S. 406, Zeile 11---15, die ergänzt wurden. \textsc{Gauss} hat hier nur die Breitenstörungen berechnet. Von der Anwendung auf die Störungen des Radiusvektors und der Länge fand sich nur der allererste Anfang. \textsc{Gauss} ging dann zur Bearbeitung der Pallasstörungen über. Die Ergebnisse finden sich im Brief an \textsc{Olbers} vom 2. Juli 1805. Sie wurden damals nicht abgedruckt.

Im Ganzen tragen \textsc{Gauss}’ Arbeiten über Ceres den Charakter von Vorarbeiten zu seinen ausgedehnten Untersuchungen über Pallas. Das erklärt, warum manche Punkte dort ausführlicher bearbeitet sind als hier.

Man vergleiche auch die bereits in Bd. VI abgedruckten einzelnen Mitteilungen über den Fortgang von \textsc{Gauss}’ Untersuchungen. An mehreren Stellen wurden kleinere Unrichtigkeiten berichtigt.
\par\nobreak\smallskip
{\raggedleft\textsc{Brendel}.\par}
\nobreak\smallskip
{\centering\rule{.38\linewidth}{.4pt}\par}
\endgroup
